% GENERATED FILE. Edit canonical lessons, then run npm run generate:books.
% canonical-source-hash: fnv1a64:07338d9168398001
% canonical-lessons: KA-C133-kappe, KA-C133-batukoli, KA-C133-handi, KA-C133-katte, KA-C133-onte

\chapter{Frog, Duck, Pig, Donkey, Camel}
\label{ch:ka-frog-duck-pig-donkey-camel}
\hlchaptermodality{133}

\begin{chapteropening}
\textbf{By the end of this chapter:} \emph{I can say a frog, a duck, a pig, a donkey and a camel.}

Complete the last lesson of chapter 133: i can say a frog, a duck, a pig, a donkey and a camel.
\end{chapteropening}

\section[kappe]{\kn{ಕಪ್ಪೆ} (kappe) — a frog}

\label{lesson:KA-C133-kappe}



\begin{quote}
Before the new one: say the Kannada for a mosquito, then the Kannada for a spider.
\end{quote}

\subsection*{You'll want to know: \kn{ಕಪ್ಪೆ}}
\textbf{\kn{ಕಪ್ಪೆ}} — \emph{kappe} — \textquotedblleft{}a frog\textquotedblright{}.

\begin{grammarlens}[title={What you've built}]
One more word for this chapter.
\end{grammarlens}

\subsection*{Your turn}
\begin{itemize}
\raggedright
  \item \emph{Say it:} \emph{kappe}
  \item \emph{Say it:} \emph{kappe}, once more
  \item \emph{Say it:} say \emph{kappe}
  \item \emph{Recall:} say the Kannada for a mosquito, then the Kannada for a spider, then say \emph{kappe} again
\end{itemize}

\begin{tcolorbox}[breakable,colback=teal!4,colframe=teal!35!black,title={Before you move on}]
What does \kn{ಕಪ್ಪೆ} mean? (\textquotedblleft{}A frog\textquotedblright{}.) Say it once more.
\end{tcolorbox}
\section[bātukōḷi]{\kn{ಬಾತುಕೋಳಿ} (bātukōḷi) — a duck}

\label{lesson:KA-C133-batukoli}



\begin{quote}
Before the new one: say the Kannada for a spider, then the Kannada for a frog.
\end{quote}

\subsection*{You'll want to know: \kn{ಬಾತುಕೋಳಿ}}
\textbf{\kn{ಬಾತುಕೋಳಿ}} — \emph{bātukōḷi} — \textquotedblleft{}a duck\textquotedblright{}.

\begin{grammarlens}[title={What you've built}]
One more word for this chapter.
\end{grammarlens}

\subsection*{Your turn}
\begin{itemize}
\raggedright
  \item \emph{Say it:} \emph{bātukōḷi}
  \item \emph{Say it:} \emph{bātukōḷi}, once more
  \item \emph{Say it:} say \emph{bātukōḷi}
  \item \emph{Recall:} say the Kannada for a spider, then the Kannada for a frog, then say \emph{bātukōḷi} again
\end{itemize}

\begin{tcolorbox}[breakable,colback=teal!4,colframe=teal!35!black,title={Before you move on}]
What does \kn{ಬಾತುಕೋಳಿ} mean? (\textquotedblleft{}A duck\textquotedblright{}.) Say it once more.
\end{tcolorbox}
\section[handi]{\kn{ಹಂದಿ} (handi) — a pig}

\label{lesson:KA-C133-handi}



\begin{quote}
Before the new one: say the Kannada for a frog, then the Kannada for a duck.
\end{quote}

\subsection*{You'll want to know: \kn{ಹಂದಿ}}
\textbf{\kn{ಹಂದಿ}} — \emph{handi} — \textquotedblleft{}a pig\textquotedblright{}.

\begin{grammarlens}[title={What you've built}]
One more word for this chapter.
\end{grammarlens}

\subsection*{Your turn}
\begin{itemize}
\raggedright
  \item \emph{Say it:} \emph{handi}
  \item \emph{Say it:} \emph{handi}, once more
  \item \emph{Say it:} say \emph{handi}
  \item \emph{Recall:} say the Kannada for a frog, then the Kannada for a duck, then say \emph{handi} again
\end{itemize}

\begin{tcolorbox}[breakable,colback=teal!4,colframe=teal!35!black,title={Before you move on}]
What does \kn{ಹಂದಿ} mean? (\textquotedblleft{}A pig\textquotedblright{}.) Say it once more.
\end{tcolorbox}
\section[katte]{\kn{ಕತ್ತೆ} (katte) — a donkey}

\label{lesson:KA-C133-katte}



\begin{quote}
Before the new one: say the Kannada for a duck, then the Kannada for a pig.
\end{quote}

\subsection*{You'll want to know: \kn{ಕತ್ತೆ}}
\textbf{\kn{ಕತ್ತೆ}} — \emph{katte} — \textquotedblleft{}a donkey\textquotedblright{}.

\begin{grammarlens}[title={What you've built}]
One more word for this chapter.
\end{grammarlens}

\subsection*{Your turn}
\begin{itemize}
\raggedright
  \item \emph{Say it:} \emph{katte}
  \item \emph{Say it:} \emph{katte}, once more
  \item \emph{Say it:} say \emph{katte}
  \item \emph{Recall:} say the Kannada for a duck, then the Kannada for a pig, then say \emph{katte} again
\end{itemize}

\begin{tcolorbox}[breakable,colback=teal!4,colframe=teal!35!black,title={Before you move on}]
What does \kn{ಕತ್ತೆ} mean? (\textquotedblleft{}A donkey\textquotedblright{}.) Say it once more.
\end{tcolorbox}
\section[oṇṭe]{\kn{ಒಂಟೆ} (oṇṭe) — a camel}

\label{lesson:KA-C133-onte}



\begin{quote}
Before the new one: say the Kannada for a pig, then the Kannada for a donkey.
\end{quote}

\subsection*{You'll want to know: \kn{ಒಂಟೆ}}
\textbf{\kn{ಒಂಟೆ}} — \emph{oṇṭe} — \textquotedblleft{}a camel\textquotedblright{}.

\begin{grammarlens}[title={What you've built}]
One more word for this chapter.
\end{grammarlens}

\subsection*{Your turn}
\begin{itemize}
\raggedright
  \item \emph{Say it:} \emph{oṇṭe}
  \item \emph{Say it:} \emph{oṇṭe}, once more
  \item \emph{Say it:} say \emph{oṇṭe}
  \item \emph{Recall:} say the Kannada for a pig, then the Kannada for a donkey, then say \emph{oṇṭe} again
\end{itemize}

\begin{tcolorbox}[breakable,colback=teal!4,colframe=teal!35!black,title={Before you move on}]
What does \kn{ಒಂಟೆ} mean? (\textquotedblleft{}A camel\textquotedblright{}.) Say it once more.
\end{tcolorbox}
